% from ega4-15-fr.tex
\oldpage[IV-3]{233}
propre (2.7.1, (vii))~; on peut donc supposer que $n(y)$ et $n(y')$ sont les nombres de points des fibres $f^{-1}(y)$ et $f'^{-1}(y')$ respectivement. Avec les notations de (i), on a alors par (15.5.2.3) et (15.5.2.1)
\begin{equation}
\tag{15.5.2.4}
\label{IV.15.5.2.4-fr}
  n(y)\leqslant\sum_i n_i(y)\leqslant\sum_i n_i(y')=n(y').
\end{equation}
et, pour que les membres extrêmes soient égaux, il faut que $n_i(y)=n_i(y')=1$ pour tout $i$. Comme on l'a vu dans (i), si $n_i(y)=1$ pour tout $i$, $X_i\cap f^{-1}(y)$ n'est pas vide et la restriction $X_i\to Y$ de $f$ est un isomorphisme~; comme les $X_i$ sont des sous-préschémas fermés de $X$, cela entraîne que $f$ est propre (\textbf{II}, 5.4.5).

(iii) On peut se borner au cas où $Y$ est affine et réduit, et comme $Y$ est par hypothèse localement intègre (\textbf{I}, 5.1.4), on peut supposer $Y$ intègre, et la fonction $y\mapsto n(y)$ constante dans $Y$. Soient encore $X_i$ ($1\leqslant i\leqslant m$) les composantes irréductibles de $X$. Il résulte de (1.10.4) que si $y'$ est le point générique de $Y$, chacun des $X_i\cap f^{-1}(y')$ est réduit à un seul point~; la restriction $f_i:X_i\to Y$ de $f$ étant un morphisme quasi-fini et dominant et tout point de $Y$ étant géométriquement unibranche, il résulte du critère de Chevalley (14.4.4) que $f_i$ est universellement ouvert. Si alors $y$ est un point quelconque de $Y$, on a $n_i(y)\leqslant n_i(y')$~; d'autre part, comme les $X_i\cap f^{-1}(y)$ sont deux à deux disjoints, on a $\sum_i n_i(y')=n(y')$~; on a donc encore les relations (15.5.2.4). Mais par hypothèse les membres extrêmes sont égaux, donc $n(y)=\sum_i n_i(y)$, ce qui implique que les $X_i\cap f^{-1}(y)$ sont deux à deux disjoints, et achève la démonstration.
\end{proof}

Combinant cette proposition avec le théorème de connexion de Zariski et ses conséquences (\textbf{III}, 4.3), on obtient les résultats suivants qui complètent (\textbf{III}, 4.3.7 et 4.3.10)~:

\begin{proposition}[15.5.3]
\label{IV.15.5.3-fr}
Soient $Y$ un préschéma irréductible localement noethérien, $f:X\to Y$ un morphisme propre~; pour tout $y\in Y$, soit $n(y)$ le nombre géométrique de composantes connexes (4.5.2) de $f^{-1}(y)$. On suppose que la restriction de $f$ à toute composante irréductible de $X$ soit un morphisme dominant dans $Y$. Alors, si $y_0$ est un point géométriquement unibranche de $Y$, la fonction $y\mapsto n(y)$ est semi-continue inférieurement au point $y_0$.
\end{proposition}

\begin{proof}
Considérons la factorisation de Stein $f=g\circ f'$ de $f$ (\textbf{III}, 4.3.3), où $g:Y'\to Y$ est un morphisme fini et $f':X\to Y$ est un morphisme surjectif dont les fibres sont géométriquement connexes (\textbf{III}, 4.3.4). Puisque $f'$ est surjectif, l'image réciproque par $f'$ de toute composante irréductible de $Y'$ contient une composante irréductible de $X$ au moins~; donc chacune des composantes irréductibles de $Y'$ domine $Y$. On peut donc appliquer le critère de Chevalley (14.4.4) aux restrictions de $g$ à chacune de ces composantes irréductibles, et on en conclut que $g$ est universellement ouvert en tout point de $g^{-1}(y_0)$. La conclusion résulte donc de (15.5.2) et du fait que le nombre géométrique des composantes connexes de $f^{-1}(y)$ est égal au nombre géométrique des points de $g^{-1}(y)$ (\textbf{III}, 4.3.4).
\end{proof}

\begin{corollary}[15.5.4]
\label{IV.15.5.4-fr}
Soient $Y$ un préschéma localement noethérien, $f:X\to Y$ un morphisme propre~; soit $y$ un point de $Y$ tel que $f$ soit universellement ouvert aux points de $f^{-1}(y)$~;
